

We now proceed to show detailed results from our ablation testing.

% \textbf{Output and CoS: } 

\begin{table}[ht]
\caption{Output and Chain-of-States Ablations}
\label{tab:abl-cos}
\centering
\small
\begin{tabular}{ll|cccc}
\toprule
Output Format & CoS & Improvement& Nonempty Improve.& Accuracy & Improved Acc.\\
\midrule
\texttt{string} & True & 7.53 & 16.12 & 46.72\% & 16.79\% \\
\texttt{string}&False & 7.03 & 19.67 & 35.77\% & 15.33\% \\
\texttt{string list}&\textbf{True} & \textbf{8.04} & \textbf{12.38} & \textbf{64.96\%} & \textbf{21.17\%} \\
\texttt{string list}&False & 7.04 & 13.58 & 51.82\% & 18.98\% \\
\texttt{string tree}&True & 7.62 & 15.34 & 49.64\% & 18.25\% \\
\texttt{string tree}&False & 6.31 & 14.17 & 44.53\% & 16.06\% \\
\bottomrule
\end{tabular}
\end{table}




By \autoref{tab:abl-cos}, we see that the optimal combination in this testing group is a \texttt{string list} output format with CoS enabled. Fix these values for all future tests.




% \textbf{Example Retrieval: }


\begin{table}[h]
\caption{Example Retrieval Ablations}
\label{tab:abl-ex}
\centering
\small
\begin{tabular}{l|cccc}
\toprule
Examples & Improvement& Nonempty Improve.& Accuracy & Improved Acc.\\
\midrule
0 &  5.67 & 8.44 & 67.15\% & 16.79\% \\
3 &  8.49 & 13.68 & 62.04\% & 19.71\% \\
5 &  8.38 & 12.9 & 64.96\% & 21.17\% \\
7 &  7.56 & 12.04 & 62.77\% & 19.71\% \\
\textbf{10} &  \textbf{9.34} & \textbf{14.7} & \textbf{63.5\%} & \textbf{21.9\%} \\
\bottomrule
\end{tabular}
\end{table}


 With the previous optimal parameters fixed, run the ablation on the number of examples. By \autoref{tab:abl-ex}, we see that the optimal combination in this testing group is $10$ examples. Fix this value for all future tests.





\begin{table}[h]
\caption{Sampling Method Ablations}
\label{tab:abl-samp}
\centering
\small
\begin{tabular}{lll|cccc}
\toprule
Method & Forward & Keep Best & Improvement& Nonempty Improve.& Accuracy & Improved Acc.\\
\midrule
None &N/A & N/A&  9.34 & 14.7 & 63.5\% & 21.9\% \\
refinement&1&False &  14.76 & 30.63 & 48.18\% & 30.66\% \\
refinement&5 &False &  12.5 & 20.88 & 59.85\% & 30.66\% \\
refinement&1 &True &  14.95 & 14.95 & 100.0\% & 30.66\% \\
refinement&5 &True &  13.15 & 13.15 & 100.0\% & 29.93\% \\
\textbf{best-of-n} &N/A&N/A&  \textbf{15.35} & \textbf{18.44} & \textbf{83.21\%} & \textbf{36.5\%} \\
\bottomrule
\end{tabular}
\end{table}


Note that forward and keep-best values are parameters for refinement of how many previous iterations to forward, and whether to keep the most recent or the best iteration in subsequent refinement steps. 



Now, with the previous optimal parameters fixed, run the ablation on the sample method. By \autoref{tab:abl-samp}, we see that the optimal combination in this testing group is best-of-n. Fix this value for all future tests.




\begin{table}[h]
\caption{Model and $n$ Ablations}
\label{tab:abl-n}
\centering
\small
\begin{tabular}{ll|cccc}
\toprule
Model & $n$ & Improvement& Nonempty Improve.& Accuracy & Improved Acc.\\
\midrule
gpt-4o & 3 &  19.66 & 24.36 & 80.7\% & 38.6\% \\
gpt-4o & 5 &  20.12 & 24.97 & 80.56\% & 36.11\% \\
gpt-4o & 7 &  22.44 & 27.21 & 82.46\% & 42.11\% \\
gpt-4o & 10 &  21.73 & 25.28 & 85.96\% & 40.35\% \\
\textbf{gpt-4o} & \textbf{15} &  \textbf{23.51} & \textbf{26.28} & \textbf{89.47\%} & \textbf{45.61\%} \\
gpt-4o-mini & 3 &  3.65 & 4.63 & 78.95\% & 8.77\% \\
gpt-4o-mini & 5 &  5.12 & 6.21 & 82.46\% & 10.53\% \\
gpt-4o-mini & 7 &  3.65 & 4.34 & 84.21\% & 8.77\% \\
gpt-4o-mini & 10 &  4.99 & 5.69 & 87.72\% & 12.28\% \\
gpt-4o-mini & 15 &  4.35 & 5.06 & 85.96\% & 12.28\% \\
gpt-4o-mini & 20 &  4.87 & 5.56 & 87.72\% & 14.04\% \\
\bottomrule
\end{tabular}
\end{table}


With the previous optimal parameters fixed, run the ablation on the value of $n$ and model. By \autoref{tab:abl-n}, we see that the optimal combination in this testing group is GPT-4o with $n=15$. Fix this value for all future tests.


\begin{table}[h]
\caption{RAG and Combination Sampling Method Ablations}
\label{tab:abl-rag}
\centering
\small
\begin{tabular}{llll|cccc}
\toprule
Combination & $m$ & $m'$  & RAG & Improvement& Nonempty Improve.& Accuracy & Improved Acc.\\
\midrule
best-of-n(refinement) & 3 & 5 & True &  33.78 & 33.78 & 100.0\% & 50.0\% \\
best-of-n(refinement) & 3 & 5 & False &  31.23 & 31.23 & 100.0\% & 46.88\% \\
best-of-n(refinement) & 5 & 3 & True &  31.85 & 31.85 & 100.0\% & 50.0\% \\
best-of-n(refinement) & 5 & 3& False &  31.35 & 31.35 & 100.0\% & 50.0\% \\
refinement(best-of-n) & 3 & 5 & True &  32.66 & 51.32 & 63.64\% & 48.48\% \\
refinement(best-of-n) & 3 & 5 & False &  32.88 & 50.1 & 65.62\% & 53.12\% \\
\textbf{refinement(best-of-n)} & \textbf{5} & \textbf{3} & \textbf{True} &  \textbf{34.88} & \textbf{57.56} & \textbf{60.61\%} & \textbf{54.55\%} \\
refinement(best-of-n) & 5 & 3 & False &  29.54 & 49.75 & 59.38\% & 43.75\% \\
best-of-n & N/A & 15 & True &  29.64 & 32.71 & 90.62\% & 56.25\% \\
best-of-n & N/A & 15 & False &  28.25 & 33.48 & 84.38\% & 53.12\% \\
\bottomrule
\end{tabular}
\end{table}


With the previous optimal parameters fixed, run the ablation on the combination methods and if RAG is enabled. By \autoref{tab:abl-rag}, we see that the optimal combination in this testing group is a $5$-step refinement with each iteration being a best-of-3 call, with RAG enabled.



